% Replaces the section currently titled
%   "Q4: obstructions to comparing the selection rules"
% Requires: \label{sec:classfunction} (Proposition~\ref{prop:coset}),
%           \label{sec:budget}, and the generated tables
%           gen_pool_equivalence.tex, gen_k_by_arm.tex.

\section{Q4: comparison of the selection rules}
\label{sec:q4}

\subsection{Statement of the question}

The question posed in the introduction conflates two claims that admit answers
of very different strength, and they are separated here.

\begin{description}
\item[Q4(i)] Does the polynomial candidate library of
  Eq.~\eqref{eq:Cpoly} select the same pool as an exhaustive search over
  the full candidate space?
\item[Q4(ii)] Does the choice of selection rule change the coupled dimension
  $K$?
\end{description}

Q4(i) is answered affirmatively, and Proposition~\ref{prop:coset} shows why the
answer is forced rather than incidental. Q4(ii) is answered only in part: a
difference is observed for one rule, but at a single orbital start it is not
separable from optimiser scatter, and it is reported as an observation.

\subsection{Design}

Three rules are compared at fixed budget, cost function, molecule and geometry:

\begin{description}
\item[exhaustive] greedy ranking over the full candidate space, i.e.\ the
  minimum-weight independent set of the matroid of
  Section~\ref{sec:classfunction}(i);
\item[iterative] greedy ranking restricted to the polynomial library
  $\Ccal^{(t)}$, re-formed at each macroiteration;
\item[quota] greedy ranking constrained to a prescribed split of single-orbital
  and quartet supports, $(n_{\mathrm{s}},n_{\mathrm{q}})=(2,\,2)$ for
  $\mathrm{H_2O}$ and $(3,\,3)$ for $\mathrm{N_2}$ at $M=\dim\Qcal$, and $(2,\,1)$ and $(3,\,2)$
  respectively at $M=\dim\Qcal-1$.
\end{description}

\paragraph{Choice of the quota split.}
The split $(n_{\mathrm{s}},n_{\mathrm{q}})$ is a prescribed input, not a
quantity the procedure derives, and it must satisfy
$n_{\mathrm{s}}+n_{\mathrm{q}}=M$. At the full budget the balanced splits
$(2,\,2)$ and $(3,\,3)$ were adopted. At the reduced budget the deficit of one was
taken from the quartets, giving $(2,\,1)$ and $(3,\,2)$, so that the number of
single-orbital rows is held fixed across the two budgets and the two campaigns
differ in one respect only. The complementary choices $(1,\,2)$ and $(2,\,3)$ were
not run. Consequently the results below characterise \emph{these} splits and
not quota selection in general: a different split could in principle come
closer to, or reach, the unconstrained optimum. This is a limitation of the
present campaign rather than a property of quota rules, and it is repeated in
Section~\ref{sec:q4limits}.

\paragraph{Sign convention.}
\begin{sloppypar}
Throughout this section
\begin{equation}
\Delta K \;=\; K_{\mathrm{quota}}-K_{\mathrm{exhaustive}},
\label{eq:dK}
\end{equation}
\noindent evaluated at fixed molecule, geometry, cost function and budget, so
that a \emph{negative} $\Delta K$ means the quota rule required a smaller coupled
space than the exhaustive rule. Comparisons between the iterative and
exhaustive rules are reported as agreement or disagreement rather than as a
signed difference, since they never disagree.
\end{sloppypar}

\paragraph{Counting of points.}
The two campaigns comprise $252$ runs: $21$ geometries ($10$ for
$\mathrm{H_2O}$, $11$ for $\mathrm{N_2}$) $\times$ $2$ cost functions $\times$
$3$ rules $\times$ $2$ budgets. A \emph{pairwise comparison} of two rules
therefore has $21\times2\times2=84$ points, and a statement restricted to one
molecule has $10\times2\times2=40$ points for $\mathrm{H_2O}$ and
$11\times2\times2=44$ for $\mathrm{N_2}$.

Each rule is followed by the same orbital optimisation of
Eq.~\eqref{eq:objective}, started from the identity, and by the same Clifford
reduction and coupled-dimension measurement, Eq.~\eqref{eq:K}. The grids are
ten $\mathrm{H_2O}$ geometries
($R=0.958$--$2.5$\,\AA) and eleven $\mathrm{N_2}$ geometries ($R=1.2$--$2.2$\,\AA),
each with both cost functions and all three rules, at two budgets: $126$
points per budget, $252$ in total. Every point satisfies the preconditions of
Section~\ref{sec:verification} --- retained reference weight $W=1$ to within
$2.2\times10^{-7}$, $n_{\mathrm{exact}}=r$, and $M$ at the recorded budget ---
so no point is excluded from what follows.

Because the rows reported by a rule are defined only up to $\Ecal$
(Section~\ref{sec:classfunction}), rules are compared by the spanned subspace
$\operatorname{span}\{\bar g_1,\dots,\bar g_M\}\subseteq\Qcal$ and by the total
selection cost $\sum_j \bar c(\bar g_j)$, never by row equality.

\subsection{Q4(i): the polynomial library attains the exhaustive optimum}

\begin{table}[htbp]
\centering
\begin{tabular}{llrrrr}
\hline
budget & cost & points & $\max|\Delta c|$ & span agree & $K$ agree\\
\hline
$M=n-\rsp$ & NC & 21 & $2.4\times10^{-14}$ & 21/21 & 21/21\\
$M=n-\rsp$ & variance & 21 & $3.0\times10^{-15}$ & 21/21 & 21/21\\
$M=n-\rsp-1$ & NC & 21 & $2.7\times10^{-14}$ & 21/21 & 21/21\\
$M=n-\rsp-1$ & variance & 21 & $2.2\times10^{-15}$ & 21/21 & 21/21\\
\hline
\end{tabular}

\caption{Agreement between the polynomial (iterative) and exhaustive rules.
``points'' counts geometries at the stated budget and cost function;
$\Delta c=\sum_j\bar c(\bar g_j)$ differs between the two rules;
``span agree'' counts points at which
$\operatorname{span}\{\bar g_j\}$ coincides; ``$K$ agree'' counts points at
which the coupled dimension coincides. Row sets never coincide.}
\label{tab:poolequiv}
\end{table}

Table~\ref{tab:poolequiv} reports agreement in three distinct senses, and the
distinction matters. Across both budgets and both cost functions the two rules
agree in \emph{total selection cost} to floating-point round-off, in the
\emph{spanned subspace} of $\Qcal$ at every point, and in the \emph{coupled
dimension} $K$ at every point. They do \emph{not} agree in the reported rows:
at no point of either campaign do the two rules return the same row set. Where
this section says the rules are ``identical'', it therefore means identical in
$\bar c$, in span, and in $K$ --- never in rows.

That combination is exactly what Section~\ref{sec:classfunction}(ii) requires:
minimum-weight independent sets of equal total weight need not coincide as
sets, and the residual freedom is the choice of coset representative, which by
Proposition~\ref{prop:coset} is unobservable. The agreement is therefore not
evidence that the polynomial library happens to be adequate; it is a
consequence of greedy optimality on the matroid, and the measurement confirms
that the implementation realises it.

The practical consequence is that the polynomial library is not an
approximation to the exhaustive search at these system sizes: it attains the
same optimum at polynomial rather than exponential candidate cost.

\subsection{Q4(ii): the effect of the selection rule on \texorpdfstring{$K$}{K}}

% h2o NC M=n-\rsp: lower 0, equal 10, higher 0, mean dK = +0.00
% h2o variance M=n-\rsp: lower 0, equal 10, higher 0, mean dK = +0.00
% n2 NC M=n-\rsp: lower 5, equal 4, higher 2, mean dK = -0.27
% n2 variance M=n-\rsp: lower 1, equal 10, higher 0, mean dK = -0.09
% h2o NC M=n-\rsp-1: lower 0, equal 10, higher 0, mean dK = +0.00
% h2o variance M=n-\rsp-1: lower 0, equal 10, higher 0, mean dK = +0.00
% n2 NC M=n-\rsp-1: lower 5, equal 5, higher 1, mean dK = -0.45
% n2 variance M=n-\rsp-1: lower 1, equal 10, higher 0, mean dK = -0.09
\begin{longtable}{lrrr}
\caption{Coupled dimension $K$ by selection rule, geometry, cost function and budget. Each block is headed by its molecule, cost function and budget. ``quota'' is the constrained rule with the splits of Section~\ref{sec:q4}; $\Delta K$ counts are in Table~\ref{tab:dk}.}
\label{tab:kbyarm}\\
\hline
$R$ (\AA) & exhaustive & iterative & quota\\
\hline
\endfirsthead
\multicolumn{4}{l}{\emph{Table \thetable\ continued}}\\
\hline
$R$ (\AA) & exhaustive & iterative & quota\\
\hline
\endhead
\hline
\endfoot
\multicolumn{4}{l}{}\\
\multicolumn{4}{l}{\emph{$\mathrm{H_2O}$, NC cost, $M=n-\rsp$}}\\
\hline
0.9580 & 4 & 4 & 4\\
1.1293 & 5 & 5 & 5\\
1.3007 & 7 & 7 & 7\\
1.4720 & 8 & 8 & 8\\
1.6433 & 8 & 8 & 8\\
1.8147 & 7 & 7 & 7\\
1.9860 & 5 & 5 & 5\\
2.1573 & 5 & 5 & 5\\
2.3287 & 4 & 4 & 4\\
2.5000 & 5 & 5 & 5\\
\multicolumn{4}{l}{}\\
\multicolumn{4}{l}{\emph{$\mathrm{H_2O}$, variance cost, $M=n-\rsp$}}\\
\hline
0.9580 & 3 & 3 & 3\\
1.1293 & 5 & 5 & 5\\
1.3007 & 6 & 6 & 6\\
1.4720 & 7 & 7 & 7\\
1.6433 & 7 & 7 & 7\\
1.8147 & 6 & 6 & 6\\
1.9860 & 5 & 5 & 5\\
2.1573 & 5 & 5 & 5\\
2.3287 & 4 & 4 & 4\\
2.5000 & 5 & 5 & 5\\
\multicolumn{4}{l}{}\\
\multicolumn{4}{l}{\emph{$\mathrm{N_2}$, NC cost, $M=n-\rsp$}}\\
\hline
1.2000 & 23 & 23 & 26\\
1.3000 & 26 & 26 & 25\\
1.4000 & 26 & 26 & 26\\
1.5000 & 26 & 26 & 27\\
1.6000 & 24 & 24 & 22\\
1.7000 & 23 & 23 & 21\\
1.8000 & 18 & 18 & 17\\
1.9000 & 14 & 14 & 13\\
2.0000 & 9 & 9 & 9\\
2.1000 & 7 & 7 & 7\\
2.2000 & 4 & 4 & 4\\
\multicolumn{4}{l}{}\\
\multicolumn{4}{l}{\emph{$\mathrm{N_2}$, variance cost, $M=n-\rsp$}}\\
\hline
1.2000 & 20 & 20 & 20\\
1.3000 & 22 & 22 & 21\\
1.4000 & 22 & 22 & 22\\
1.5000 & 23 & 23 & 23\\
1.6000 & 22 & 22 & 22\\
1.7000 & 21 & 21 & 21\\
1.8000 & 17 & 17 & 17\\
1.9000 & 13 & 13 & 13\\
2.0000 & 9 & 9 & 9\\
2.1000 & 7 & 7 & 7\\
2.2000 & 4 & 4 & 4\\
\multicolumn{4}{l}{}\\
\multicolumn{4}{l}{\emph{$\mathrm{H_2O}$, NC cost, $M=n-\rsp-1$}}\\
\hline
0.9580 & 1 & 1 & 1\\
1.1293 & 1 & 1 & 1\\
1.3007 & 1 & 1 & 1\\
1.4720 & 1 & 1 & 1\\
1.6433 & 1 & 1 & 1\\
1.8147 & 1 & 1 & 1\\
1.9860 & 1 & 1 & 1\\
2.1573 & 1 & 1 & 1\\
2.3287 & 1 & 1 & 1\\
2.5000 & 3 & 3 & 3\\
\multicolumn{4}{l}{}\\
\multicolumn{4}{l}{\emph{$\mathrm{H_2O}$, variance cost, $M=n-\rsp-1$}}\\
\hline
0.9580 & 1 & 1 & 1\\
1.1293 & 1 & 1 & 1\\
1.3007 & 1 & 1 & 1\\
1.4720 & 1 & 1 & 1\\
1.6433 & 1 & 1 & 1\\
1.8147 & 1 & 1 & 1\\
1.9860 & 1 & 1 & 1\\
2.1573 & 1 & 1 & 1\\
2.3287 & 1 & 1 & 1\\
2.5000 & 3 & 3 & 3\\
\multicolumn{4}{l}{}\\
\multicolumn{4}{l}{\emph{$\mathrm{N_2}$, NC cost, $M=n-\rsp-1$}}\\
\hline
1.2000 & 16 & 16 & 15\\
1.3000 & 18 & 18 & 17\\
1.4000 & 18 & 18 & 19\\
1.5000 & 18 & 18 & 18\\
1.6000 & 17 & 17 & 16\\
1.7000 & 25 & 25 & 23\\
1.8000 & 19 & 19 & 19\\
1.9000 & 15 & 15 & 14\\
2.0000 & 7 & 7 & 7\\
2.1000 & 4 & 4 & 4\\
2.2000 & 3 & 3 & 3\\
\multicolumn{4}{l}{}\\
\multicolumn{4}{l}{\emph{$\mathrm{N_2}$, variance cost, $M=n-\rsp-1$}}\\
\hline
1.2000 & 11 & 11 & 11\\
1.3000 & 15 & 15 & 14\\
1.4000 & 17 & 17 & 17\\
1.5000 & 16 & 16 & 16\\
1.6000 & 15 & 15 & 15\\
1.7000 & 15 & 15 & 15\\
1.8000 & 11 & 11 & 11\\
1.9000 & 7 & 7 & 7\\
2.0000 & 5 & 5 & 5\\
2.1000 & 4 & 4 & 4\\
2.2000 & 2 & 2 & 2\\
\end{longtable}


Three observations follow, in decreasing order of evidential strength.

\paragraph{Cost-ranked rules do not differ.}
The exhaustive and iterative columns of Table~\ref{tab:kbyarm} carry the same
$K$ at all $84$ points, in the sense fixed above (equal $\bar c$, equal span,
equal $K$; unequal rows). By the argument of
Section~\ref{sec:classfunction}(ii) this is structural, and it is therefore not
an independent comparison of two rules: they realise the same optimum.

\paragraph{The quota rule as configured is suboptimal in selection cost.}
The quota constraint prevents the greedy optimum from being reached, and the
recorded total selection cost of the quota pool exceeds that of the exhaustive
pool at every one of the $84$ points, never falling below it. The constraint is
thus active at every geometry, and the quota arm is a genuinely different
selection rather than a relabelling. This is a statement about the splits
listed in Section~\ref{sec:q4} only; a different admissible split
$(n_{\mathrm{s}},n_{\mathrm{q}})$ with $n_{\mathrm{s}}+n_{\mathrm{q}}=M$ could
in principle attain the unconstrained optimum, and no such split was tested.

\paragraph{A lower selection cost did not produce a lower $K$.}
Despite being worse under the objective that the ranking minimises, the quota
rule did not require a larger coupled space. Table~\ref{tab:dk} counts
$\Delta K$ of Eq.~\eqref{eq:dK} as (lower, equal, higher) over the geometries
of each block.

\begin{table}[htbp]
\centering
\begin{tabular}{llll}
\hline
system & cost & $M=n-\rsp$ & $M=n-\rsp-1$\\
\hline
$\mathrm{H_2O}$ & NC       & $(0,10,0)$, $\overline{\Delta K}=+0.00$ & $(0,10,0)$, $+0.00$\\
$\mathrm{H_2O}$ & variance & $(0,10,0)$, $+0.00$ & $(0,10,0)$, $+0.00$\\
$\mathrm{N_2}$  & NC       & $(5,4,2)$, $-0.27$ & $(5,5,1)$, $-0.45$\\
$\mathrm{N_2}$  & variance & $(1,10,0)$, $-0.09$ & $(1,10,0)$, $-0.09$\\
\hline
\end{tabular}
\caption{Quota rule against the exhaustive rule, counted as
(lower, equal, higher) in $\Delta K=K_{\mathrm{quota}}-K_{\mathrm{exhaustive}}$
over the $10$ ($\mathrm{H_2O}$) or $11$ ($\mathrm{N_2}$) geometries of each
block, with the mean $\overline{\Delta K}$. Negative favours the quota rule.}
\label{tab:dk}
\end{table}

The effect is confined to $\mathrm{N_2}$ and predominantly to the NC cost. For
$\mathrm{H_2O}$ the quota pool is more expensive in $\bar c$ at all $40$ points
($10$ geometries $\times$ $2$ cost functions $\times$ $2$ budgets) yet returns
exactly the same $K$ at all $40$; under the parity variance the $\mathrm{N_2}$
difference is a single geometry at each budget. Where a difference does appear
its sign is not consistent: over $\mathrm{N_2}$ and both cost functions the
range of $\Delta K$ is $[-2,+3]$ at $M=n-\rsp$ and $[-2,+1]$ at $M=n-\rsp-1$,
the largest single departure being $+3$ at $R=1.2$\,\AA.

This bears on the interpretation of the NC score. The exhaustive rule minimises
$\bar c$ exactly, whereas $K$ is measured after the orbital optimisation of
Eq.~\eqref{eq:objective}; the two are separated by that optimisation. A rule
that is worse in $\bar c$ and no worse in $K$ is therefore not a contradiction.
What these data establish is narrow and is stated as such: \emph{in the present
configuration} --- two molecules in a minimal basis, the quota splits listed
above, one orbital start --- ordering candidate rows by $\bar c$ did not order
the resulting $K$, and at $40$ of the points a strictly worse $\bar c$ produced
an identical $K$. Whether $\bar c$ is monotonically related to $K$ in general
is not addressed here, and no such claim is made; the observation is that the
relationship failed to hold in this configuration. In particular the exhaustive
column must not be read as a lower bound on $K$: it minimises the selection
cost, not the coupled dimension.

\subsection{What these data do not establish}
\label{sec:q4limits}

\begin{enumerate}
\item \textbf{No scatter estimate.} A single orbital start (the identity) was
run at each point, so no within-configuration variance is available. The
observed differences in $K$ are of the same order as the shifts seen between
budgets at fixed rule, and a difference of one or two units must be read as
``no detectable difference at one start''. Deciding whether the quota effect is
real requires repeating a subset with randomised starts at fixed pool; that
study is identified in Section~\ref{sec:conclusions} and has not been performed.
\item \textbf{$\mathrm{H_2O}$ is uninformative at the reduced budget.} At
$M=\dim\Qcal-1$ the decoupled energy already lies within $\varepsilon_E$ of
$E_{\mathrm{target}}$ at nine of ten $\mathrm{H_2O}$ geometries, giving $K=1$. These
are floor values, not agreement between rules, and they carry no discriminating
information; the $\mathrm{N_2}$ rows carry the entire Q4(ii) comparison.
\item \textbf{One quota split per budget.} The constrained rule was run at
$(n_{\mathrm{s}},n_{\mathrm{q}})=(2,2)/(3,3)$ and $(2,1)/(3,2)$ only. The
complementary splits were not tested, so ``the quota rule is suboptimal in
$\bar c$'' is established for these splits and not for quota selection as such.
\item \textbf{Spin-symmetric ground set.} Candidates are products of $\Zb_p$
(Section~\ref{sec:classfunction}(0)), so every optimum reported here is an
optimum over $\iota(\FF^{n})$, not over $\FF^{N}$. A spin-resolved search may
admit lower-cost pools at the same $M$.
\item \textbf{Scope.} Two molecules in a minimal basis, one orbital-rotation
parameterisation, and one reference convention. Nothing here establishes
behaviour at larger basis sets or for systems whose exact group is not a point
group supplemented by particle-number parities.
\item \textbf{Near-degenerate roots.}\begin{sloppypar}
The parent FCI is solved with \texttt{pyscf.fci.direct\_spin1} at
\texttt{conv\_tol}\,$=10^{-10}$ and \texttt{max\_cycle}\,$=10^{4}$, from
PySCF's deterministic initial guess; no
random seed enters. Where the FCI gap is small --- $\mathrm{H_2O}$ at
$R\ge2.33$\,\AA, gap $4.5\times10^{-4}$\,Ha --- no explicit tie-break is
imposed on near-degenerate eigenvectors, so the selected vector is reproducible
for a fixed build but is not guaranteed to be stable across BLAS or PySCF
versions. Energies are unaffected; overlap-ordered quantities in principle are.
\end{sloppypar}
\end{enumerate}
